\section{Introduction}\label{sec:Introduction}
% The abstract stands alone and expands acronyms on its own. Here, the body starts fresh,
% so that its first \ac expands the acronym again.
\glsresetall

Texture mapping adds detail to surfaces without adding geometry~\cite{Catmull74SAC}.
A rasterizer interpolates the texture coordinates of the vertices across each triangle.
\term{Affine texture mapping} interpolates them linearly in screen space.
It ignores the perspective and bends the texture on surfaces that recede from the viewer, as \cref{fig:Teaser}~(a) shows.
The first PlayStation mapped textures this way, and some of its games split large polygons into smaller ones to reduce the bending~\cite{Copetti19PAP}.

\term{Perspective-correct} interpolation~\cite{Heckbert91IPT,Blinn92HI} removes the error.
It needs a division per pixel, which today's \acp{GPU} perform in hardware.
Without this division, a renderer must subdivide its surfaces instead.
How finely?
This paper answers the question in closed form.
Thereby, we make the following contributions:
\begin{itemize}
    \item We derive the largest error of affine interpolation along an edge in closed form (\cref{sec:EdgeError}).
    \item We bound the error of subdivided surfaces and derive how many pieces keep it below a threshold (\cref{sec:Subdivision}).
    \item We confirm both results with a software rasterizer (\cref{sec:Results}).
\end{itemize}
We measure the error of the texture coordinates for nearest-texel lookups.
Texture filtering~\cite{Williams83PP} is beyond the scope of this paper.
